\section{Discussion}
\noindent\textbf{System overhead}.
VibVoice can transmit raw acceleration data to the mobile device at a bit rate of 153.6 kbps (6*16 bits*$1.6\,\text{kHz}$). Bluetooth profiles such as Hands-Free Profile (HFP) support a 16 kbps ($16\,\text{kHz}$) audio data rate, which is adequate for transmitting acceleration data under Bluetooth 5.0 \citep{Bluetoot0:online}. IMU data can also be compressed to further reduce communication overhead.

Most head-mounted wearables run on the battery. The extra energy consumed by VibVoice is mainly ADC, which is 0.54 mW (i.e., $180 \mu A * 3 V$) \citep{Inertial63:online}. In contrast, each AirPods Pro earbud has a 43 mAh battery, which can support 3.5 hours of talking time. Hence, VibVoice only has an extra $1.5\%$ power consumption for earphones. Note that many earphones like AirPods Pro already have motion-related applications running (e.g., spatial audio \citep{spatial-audio}), in which VibVoice can share the data collection process.

\noindent\textbf{Bone conduction function}.
VibVoice shows that Bone Conduction Function can be estimated using the Gaussian distribution with a fine-tuning process. We note that some finite-element models can be used for bone-conduction sound simulation, which further improves the performance of VibVoice \citep{chang2016development}.

\noindent\textbf{Language}.
The data augmentation approach is based on the bone conduction effect of the user's skull, which contains no language-specific features, which means VibVoice can be generalized to any language.

\noindent\textbf{Sensing location}.
Based on our experience of extensive evaluations, we summarise two guidelines about the optimal placement of the IMU sensors on the user's head for the HMW manufacturers: 1) closer to the vibrating organ, and 2) a tight contact with the head. Meanwhile, it is also important to consider human comfort and compatibility with current devices.

\noindent\textbf{Customized device}.
Future implementations could combine VibVoice with raw microphone data before hardware-level signal processing. In the current prototype, the recorded audio is processed by the device's built-in acoustic pipeline, which requires fewer system privileges and no hardware modification. However, this black-box processing can bias the input. For example, commercial earphones may suppress target speech or introduce artifacts that do not appear in the training dataset. VibVoice could perform better with access to raw microphone-array data from the TWS or HMW platform, such as the signals used for active noise cancellation or acoustic beamforming.
